% A4縦・1段組の学習用（教科書風）
\documentclass[paper=a4,fleqn]{jlreq}
\special{papersize=\the\paperwidth,\the\paperheight}

\usepackage[margin=20truemm]{geometry}
\usepackage[dvipdfmx]{graphicx}
\usepackage{amsmath,amssymb}
\usepackage{tikz}
\usepackage{pgfplots}
\usepackage{etoolbox}
\usepackage{needspace}
\usepackage{bm}
\pgfplotsset{compat=1.18}
\usetikzlibrary{arrows.meta,calc,angles,quotes}

% If there is not enough space left, start the section on the next page.
\preto{\section}{\needspace{0.35\textheight}}
\preto{\subsection}{\needspace{0.32\textheight}}

\ModifyHeading{section}{
  lines=1,
  before_space=14pt,
  after_space=6pt
}

\ModifyHeading{subsection}{
  lines=1,
  before_space=10pt,
  after_space=4pt
}

\title{ベクトル}
\author{}
\date{}

\makeatletter
\renewcommand{\maketitle}{%
  \begin{center}
    {\large\bfseries \@title\par}
  \end{center}
  \vspace{4pt}
}
\makeatother

% 公式・キーポイントを目立たせるコマンド
\newcommand{\formula}[1]{%
  \vspace{4pt}%
  \begin{center}{\large\bfseries\boldmath\(\displaystyle #1\)}\end{center}%
  \vspace{4pt}%
}

\begin{document}
\maketitle

%==========================================================
\section{ベクトルとは何か}
%==========================================================

\subsection{スカラーとベクトル}

これまで扱ってきた数値（例：8, -2, 3.5）は\textbf{スカラー}と呼ばれ、量だけを表す。

一方、\textbf{ベクトル}は\textbf{量と方向の両方}を表す。

\begin{center}
{\large\bfseries
  スカラー：量のみ（8, -5など）\\[4pt]
  ベクトル：量 + 方向（速度、位置変化など）
}
\end{center}

ベクトルは矢印で表す：\(\vec{a}\) または \(\overrightarrow{AB}\)

\subsection{ベクトルがなぜ必要か}

\textbf{高校物理では使わないが、ゲーム数学では必須。}

ゲーム開発では、物体の位置・速度・移動を計算する必要があり、すべて向きと量を持つ。ベクトルを使えば、向きと量をまとめて扱えるため、計算が簡潔かつ直感的。

\textbf{例：}
\begin{itemize}
  \item 物体の速度：x方向に毎秒 2、y方向に毎秒 3 \(\rightarrow\) \(\vec{v} = \begin{pmatrix} 2 \\ 3 \end{pmatrix}\)[m/s]
  \item 物体の移動方向：右上45°、距離 5 \(\rightarrow\) ベクトルで表現
  \item キャラの力の向き：斜め下向きに 10N \(\rightarrow\) 方向と大きさをまとめて表現
\end{itemize}

\subsection{ベクトルの大小と等しさ}

2つのベクトルが\textbf{等しい}とは、\textbf{大きさと向きが同じ}ことである。
位置は関係ない（どこから始まろうと同じベクトル）。

\textbf{例：} 線分 AB と線分 CD の長さと向きが同じなら、\(\overrightarrow{AB} = \overrightarrow{CD}\)

\subsection{確認問題}
\begin{enumerate}
  \item スカラーとベクトルの違いを説明せよ。
  \item ベクトルがゲーム開発で必要な理由を述べよ。
  \item 2つのベクトルが等しいための条件は何か。
\end{enumerate}

\newpage
%==========================================================
\section{ベクトルの成分表示}
%==========================================================

\subsection{成分とは何か}

2次元座標上で、ベクトルを\textbf{成分}で表す。

原点 O から出発したベクトル \(\vec{a}\) が点 \((a_x, a_y)\) に到達するとき：

\formula{\vec{a} = \begin{pmatrix} a_x \\ a_y \end{pmatrix}}

\textbf{重要：成分は、ベクトルが原点から出発した場合の座標である。始点がどこにあっても、成分は変わらない。}

\textbf{例：}
\begin{itemize}
  \item \(\vec{a} = \begin{pmatrix} 3 \\ 2 \end{pmatrix}\) は「右に3、上に2」
  \item \(\vec{b} = \begin{pmatrix} -1 \\ 4 \end{pmatrix}\) は「左に1、上に4」
  \item \(\vec{c} = \begin{pmatrix} 2 \\ -3 \end{pmatrix}\) は「右に2、下に3」
\end{itemize}

\subsection{成分と大きさ（ピタゴラスの定理）}

ベクトル \(\vec{a} = \begin{pmatrix} a_x \\ a_y \end{pmatrix}\) の大きさは、ピタゴラスの定理から：

\formula{|\vec{a}| = \sqrt{a_x^2 + a_y^2}}

\textbf{例：} \(\vec{a} = \begin{pmatrix} 3 \\ 4 \end{pmatrix}\) のとき、\(|\vec{a}| = \sqrt{9 + 16} = \sqrt{25} = 5\)

\subsection{点と点を結ぶベクトル}

点 \(A(x_1, y_1)\) から点 \(B(x_2, y_2)\) へのベクトルは：

\formula{\overrightarrow{AB} = \begin{pmatrix} x_2 - x_1 \\ y_2 - y_1 \end{pmatrix}}

\textbf{例：} \(A(1, 2)\)、\(B(5, 6)\) のとき、\(\overrightarrow{AB} = \begin{pmatrix} 4 \\ 4 \end{pmatrix}\)

\subsection{確認問題}
\begin{enumerate}
  \item ベクトル \(\vec{a} = \begin{pmatrix} 2 \\ -1 \end{pmatrix}\) の大きさを求めよ。
  \item 点 \(A(1, 3)\)、点 \(B(5, 1)\) のとき、\(\overrightarrow{AB}\) を成分で表せ。
  \item \(\vec{c} = \begin{pmatrix} -3 \\ 0 \end{pmatrix}\) の大きさは？
\end{enumerate}

\newpage
%==========================================================
\section{ベクトルの演算}
%==========================================================

\subsection{ベクトルの加算}

2つのベクトルを足すとき、\textbf{成分ごとに足す}。

\formula{\vec{a} + \vec{b} = \begin{pmatrix} a_x + b_x \\ a_y + b_y \end{pmatrix}}

\textbf{例：}
\begin{itemize}
  \item \(\begin{pmatrix} 2 \\ 3 \end{pmatrix} + \begin{pmatrix} 1 \\ -1 \end{pmatrix} = \begin{pmatrix} 3 \\ 2 \end{pmatrix}\)
  \item ゲーム例：速度 \(\vec{v}_1 = (2, 3)\) で移動しながら、さらに加速度 \((1, -1)\) が加わる \(\rightarrow\) 合成速度は \((3, 2)\)
\end{itemize}

\subsection{ベクトルの減算}

\formula{\vec{a} - \vec{b} = \begin{pmatrix} a_x - b_x \\ a_y - b_y \end{pmatrix}}

\textbf{例：} \(\begin{pmatrix} 5 \\ 2 \end{pmatrix} - \begin{pmatrix} 2 \\ 3 \end{pmatrix} = \begin{pmatrix} 3 \\ -1 \end{pmatrix}\)

\subsection{スカラー倍}

実数 \(k\) とベクトルを掛けると、成分がすべて \(k\) 倍になる。

\formula{k\vec{a} = \begin{pmatrix} ka_x \\ ka_y \end{pmatrix}}

\begin{itemize}
  \item \(k > 0\)：同じ向き、大きさが \(k\) 倍
  \item \(k < 0\)：向きが反対、大きさが \(|k|\) 倍
  \item \(k = 0\)：零ベクトル
\end{itemize}

\textbf{例：} \(3 \begin{pmatrix} 2 \\ -1 \end{pmatrix} = \begin{pmatrix} 6 \\ -3 \end{pmatrix}\)

\textbf{ゲーム例：}速度ベクトルを2倍にすると、物体は2倍の速さで移動する。

\subsection{確認問題}
\begin{enumerate}
  \item \(\vec{a} = \begin{pmatrix} 1 \\ 2 \end{pmatrix}\)、\(\vec{b} = \begin{pmatrix} 3 \\ -1 \end{pmatrix}\) のとき、\(\vec{a} + \vec{b}\) と \(\vec{a} - \vec{b}\) を求めよ。
  \item \(2\vec{a} + 3\vec{b}\) を求めよ。
  \item 速度ベクトル \(\vec{v} = \begin{pmatrix} 1 \\ 1 \end{pmatrix}\) で 5 秒間移動する。移動ベクトルは？
\end{enumerate}

\newpage
%==========================================================
\section{内積となす角}
%==========================================================

\subsection{内積とは}

2つのベクトルの\textbf{内積}（スカラー積）は：

\formula{\vec{a} \cdot \vec{b} = a_x b_x + a_y b_y}

\textbf{重要：内積の結果はベクトルではなく、スカラー（数値）である。}

\textbf{例：} \(\begin{pmatrix} 2 \\ 3 \end{pmatrix} \cdot \begin{pmatrix} 1 \\ 4 \end{pmatrix} = 2 \cdot 1 + 3 \cdot 4 = 2 + 12 = 14\)

\subsection{内積と角度の関係}

内積には別の定義もある：

\formula{\vec{a} \cdot \vec{b} = |\vec{a}||\vec{b}|\cos\theta}

ここで \(\theta\) は2つのベクトルのなす角。

\textbf{この2つの式から、なす角を求められる：}

\formula{\cos\theta = \frac{\vec{a} \cdot \vec{b}}{|\vec{a}||\vec{b}|}}

\subsection{垂直判定（ゲーム応用）}

\(\cos\theta = 0\) のとき \(\theta = 90°\) なので：

\textbf{\(\vec{a} \cdot \vec{b} = 0\) \(\Rightarrow\) \(\vec{a} \perp \vec{b}\)（垂直）}

ゲーム開発では、2つのベクトルが垂直かどうかを調べるのに便利。

\textbf{例：} \(\vec{a} = \begin{pmatrix} 2 \\ 1 \end{pmatrix}\)、\(\vec{b} = \begin{pmatrix} -1 \\ 2 \end{pmatrix}\) のとき、

\(\vec{a} \cdot \vec{b} = 2 \cdot (-1) + 1 \cdot 2 = -2 + 2 = 0\) \(\Rightarrow\) \textbf{垂直}

\subsection{確認問題}
\begin{enumerate}
  \item \(\vec{a} = \begin{pmatrix} 1 \\ 2 \end{pmatrix}\)、\(\vec{b} = \begin{pmatrix} 3 \\ 1 \end{pmatrix}\) のとき、\(\vec{a} \cdot \vec{b}\) を求めよ。
  \item \(\vec{a} = \begin{pmatrix} 2 \\ 1 \end{pmatrix}\)、\(\vec{b} = \begin{pmatrix} -1 \\ 2 \end{pmatrix}\) は垂直か確認せよ。
  \item \(\vec{a} = \begin{pmatrix} 1 \\ 1 \end{pmatrix}\)、\(\vec{b} = \begin{pmatrix} 2 \\ 0 \end{pmatrix}\) のなす角を求めよ（\(|\vec{a}| = \sqrt{2}\)、\(|\vec{b}| = 2\)）。
\end{enumerate}

\newpage
%==========================================================
\section{ゲーム開発への応用}
%==========================================================

\subsection{速度ベクトルと位置}

ゲーム開発では、物体の\textbf{速度をベクトル}で表す。

速度ベクトル：\(\vec{v} = \begin{pmatrix} v_x \\ v_y \end{pmatrix}\) [m/s]

初期位置：\(p_0 = (x_0, y_0)\)

\textbf{t秒後の位置：}

\formula{p(t) = p_0 + t \cdot \vec{v} = \begin{pmatrix} x_0 + t \cdot v_x \\ y_0 + t \cdot v_y \end{pmatrix}}

\textbf{例：}初期位置 \(p_0 = (2, 3)\)、速度 \(\vec{v} = (1, 2)\) [m/s]、時間 \(t = 3\) [s]

\[p(3) = (2, 3) + 3 \cdot (1, 2) = (2 + 3, 3 + 6) = (5, 9)\]

物体は 3 秒後に座標 \((5, 9)\) に到達する。

\subsection{衝突判定（円同士）}

ゲームで2つの円が衝突したかを判定する場合、\textbf{中心間の距離}を調べる。

円 A：中心 \(A(x_a, y_a)\)、半径 \(r_a\)

円 B：中心 \(B(x_b, y_b)\)、半径 \(r_b\)

\textbf{中心間のベクトル：}\(\overrightarrow{AB} = \begin{pmatrix} x_b - x_a \\ y_b - y_a \end{pmatrix}\)

\textbf{中心間の距離：}\(d = |\overrightarrow{AB}| = \sqrt{(x_b - x_a)^2 + (y_b - y_a)^2}\)

\textbf{衝突条件：}
\[
  d < r_a + r_b \quad \Rightarrow \quad \text{衝突している}
\]

\textbf{衝突判定の高速化（根号を避ける）：}
\[
  d^2 < (r_a + r_b)^2
\]

\textbf{例：}円A（中心 (1, 2)、半径 1）、円B（中心 (4, 6)、半径 2）

\(\overrightarrow{AB} = (3, 4)\)、\(d = \sqrt{9 + 16} = \sqrt{25} = 5\)

\(r_a + r_b = 1 + 2 = 3\)

\(d = 5 > 3\) \(\Rightarrow\) \textbf{衝突していない}

\subsection{確認問題}
\begin{enumerate}
  \item 初期位置 \(p_0 = (0, 0)\)、速度 \(\vec{v} = (2, -1)\) [m/s] のとき、\(t = 2\) [s]後の位置を求めよ。
  \item 円A（中心 (0, 0)、半径 1）と円B（中心 (1, 1)、半径 1）は衝突しているか？
  \item 3秒間で位置 (1, 1) から (7, 7) へ移動する物体の速度ベクトルを求めよ。
\end{enumerate}

\end{document}
